\section{Background}
\label{background}

\subsection*{Hebbian learning}
The fundamental capacity of the brain to learn and adapt does not arise from the intelligence of individual neurons, but from the dynamic plasticity of the connections between them. This principle is articulated by Donald Hebb (1949), who proposed that ``cells that fire together, wire together'' \cite{hebb2005organization}. Hebbian theory implies that synaptic strength is increased when the presynaptic cell repeatedly contributes to the firing of the postsynaptic cell, forming the cellular basis of associative learning \cite{brown2020donald, lazari2022hebbian}. The bidirectional refinement of this process is captured by long-term potentiation (LTP), in which co-active synapses are strengthened, and long-term depression (LTD), in which weakly correlated or temporally misaligned activity weakens connections \cite{lazari2022hebbian}.

More recent work extends this to reward-modulated Hebbian plasticity, where synaptic updates also depend on a global modulatory signal (e.g., dopamine) carrying reward or prediction error \cite{kusmierz2017learning}. This three-factor rule allows neural circuits to reinforce only those causal associations that lead to positive outcomes. Through this mechanism, transient neural co-activations stabilise into functional assemblies capable of complex computation.

\subsection*{Social learning}
Human intelligence is fundamentally social: the Social Brain Hypothesis posits that the evolutionary expansion of the primate brain was driven by the computational demands of managing complex, large-scale social relationships \cite{dunbar1998social}. Beyond direct experience, humans rely on social learning to acquire skills and knowledge from one another, allowing groups to outperform any individual member \cite{tomasello1993cultural, boyd2011cultural}. Crucially, social learning is selective rather than indiscriminate: it depends on persistent relationships that determine whom to learn from and what to retain \cite{kendal2018social}.

\subsection*{Dec-POMDPs and cooperative MARL}
Multi-agent environments are formally modelled as Decentralised Partially Observable Markov Decision Processes (Dec-POMDPs) \cite{bernstein2002complexity}. In a Dec-POMDP, $N$ agents act simultaneously in a shared environment under partial observability: each agent $i$ receives a local observation $o_i$, selects an action $a_i$ according to its policy $\pi_{\theta_i}(a_i \mid o_i)$, and receives a reward $r_i$. In the cooperative setting, agents share a common reward signal and the joint objective is to maximise the expected return.

Dec-POMDPs capture the key challenges of real-world multi-agent systems, including partial observability, non-stationarity induced by simultaneous learning, and the need for coordination under limited information \cite{wong2023deep}. Optimal solutions are intractable in general, so most MARL methods approximate Dec-POMDP solutions through centralised training with decentralised execution (CTDE) \cite{amato2015scalable}, in which agents condition value estimates on global information during training while restricting actor inputs to local observations at execution.

\subsection*{Proximal Policy Optimisation (PPO) and MAPPO}
PPO is a policy-gradient method that improves training stability by replacing the standard objective with a clipped surrogate that prevents large parameter updates from destabilising the learned policy \cite{yu2022surprising}. The objective constrains the probability ratio between the new and old policies, ensuring updates remain within a trust region. PPO is augmented with a value-function loss and an entropy bonus that discourages premature convergence to deterministic policies.

Multi-Agent PPO (MAPPO) extends PPO to cooperative MARL under the CTDE paradigm \cite{yu2022surprising}: each agent maintains a local actor that conditions on its observation, while a centralised critic conditions on the joint observation of all agents during training. The critic's advantage estimates account for the team's contribution rather than each agent's individual reward signal.

\subsection*{LLM-based agents}
LLMs have recently been used as the central reasoning core of autonomous agents, mapping structured natural-language observations to discrete actions through prompt-conditioned generation \cite{xi2025rise}. To allow such agents to learn from environmental reward, the LLM is typically frozen and a small set of trainable parameters is added: a classification head that maps a pooled hidden state to a distribution over the action space, and low-rank adapters (LoRA) inserted into the attention layers \cite{hu2022lora}, which constrain the policy update to a low-dimensional subspace and prevent the reward signal from corrupting the LLM's general reasoning capabilities. This paradigm enables policy optimisation against environment reward while preserving the LLM's pre-trained knowledge \cite{qian2025toolrl}.
